\section{What we don't know about recovering from methamphetamine}\label{14_what_we_dont_know-what-we-dont-know-about-recovering-from-methamphetamine}

\textbf{The central finding of this review is an absence.} Despite over 5,000 publications on amphetamine-type stimulant neurotoxicity, no study has systematically tracked whether --- or when --- former methamphetamine users regain the capacity to experience pleasure, motivation, and a meaningful life beyond roughly 14 months of abstinence. The foundational neuroimaging evidence for ``brain recovery'' rests on longitudinal samples of five people. The most clinically devastating symptom of protracted withdrawal --- anhedonia --- has never been measured with a validated instrument across the recovery timeline. No FDA-approved medication exists for methamphetamine use disorder. \href{https://www.sciencedirect.com/science/article/abs/pii/S009174352300230X}{ScienceDirect} \href{https://nida.nih.gov/publications/drugs-brains-behavior-science-addiction/treatment-recovery}{NIDA} The field has generated enormous knowledge about how the drug damages brain tissue in mice and almost none about what it is like to be a human being trying to recover. This gap is not accidental. It is the predictable product of funding structures, career incentives, species selection bias, and a measurement paradigm that equates abstinence with recovery. What follows is a systematic accounting of both what is known and what has been left unexamined.

\begin{center}\rule{0.5\linewidth}{0.5pt}\end{center}

\subsection{The neurochemical evidence rests on remarkably thin foundations}\label{14_what_we_dont_know-the-neurochemical-evidence-rests-on-remarkably-thin-foundations}

The claim that ``the brain heals in 14--18 months'' --- ubiquitous in treatment center marketing and popular recovery literature --- traces to two landmark papers, both with longitudinal samples of \textbf{n = 5}.

Volkow et al.~(2001, \emph{Journal of Neuroscience}) used PET with {[}¹¹C{]}d-threo-methylphenidate to measure dopamine transporter (DAT) density in 12 methamphetamine abusers, of whom only 5 remained abstinent long enough to be retested at 12--17 months. \href{https://pmc.ncbi.nlm.nih.gov/articles/PMC6763886/}{PubMed Central} \href{https://nida.nih.gov/publications/drugs-brains-behavior-science-addiction/treatment-recovery}{NIDA} DAT availability increased by \textbf{+19\% in caudate and +16\% in putamen} \href{https://pmc.ncbi.nlm.nih.gov/articles/PMC6763886/}{PubMed Central +2} --- significant, but still below control levels. Critically, the authors noted that ``neuropsychological tests did not improve to the same extent,'' \href{https://pubmed.ncbi.nlm.nih.gov/11717374/}{PubMed} explicitly warning against equating DAT recovery with functional restoration. \href{https://www.jneurosci.org/content/21/23/9414}{Journal of Neuroscience +2}

The companion study, Wang et al.~(2004, \emph{American Journal of Psychiatry}), measured brain metabolism with FDG-PET in the same 5 longitudinal subjects plus 8 cross-sectional participants. Thalamic metabolism recovered and correlated with improved motor and verbal memory performance. \href{https://psychiatryonline.org/doi/10.1176/appi.ajp.161.2.242}{Psychiatry Online} But \textbf{striatal metabolism --- particularly in the nucleus accumbens --- did not recover}, even after 12--17 months. \href{https://ohsu.elsevierpure.com/en/publications/partial-recovery-of-brain-metabolism-in-methamphetamine-abusers-a-2/}{Elsevierpure} The authors wrote that ``persistent decreases in the nucleus accumbens could account for the persistence of amotivation and anhedonia.'' \href{https://news.unboundmedicine.com/medline/citation/14754772/Partial_recovery_of_brain_metabolism_in_methamphetamine_abusers_after_protracted_abstinence_}{Unboundmedicine} This dissociation --- thalamic recovery enabling some cognitive improvement while the brain's reward center remains hypometabolic --- is the closest the literature comes to a neurobiological explanation for why abstinent people can pass cognitive tests while reporting they cannot enjoy life.

The Volkow 2015 paper (\emph{NeuroImage}) fundamentally reinterpreted the earlier DAT findings. In 16 methamphetamine abusers (9 retested at 9 months), DAT increased approximately 20\%, but \textbf{methylphenidate-induced dopamine release did not change}. \href{https://www.sciencedirect.com/science/article/abs/pii/S1053811915006461}{ScienceDirect} The conclusion was striking: ``The loss of DAT in these {[}methamphetamine abusers{]} does not reflect degeneration of dopamine terminals.'' \href{https://www.sciencedirect.com/science/article/abs/pii/S1053811915006461}{ScienceDirect +2} DAT ``recovery'' likely represents adaptive upregulation of a regulatory protein, not regeneration of damaged neural architecture. \href{https://pmc.ncbi.nlm.nih.gov/articles/PMC6763886/}{PubMed Central} D2/D3 receptor availability in the caudate was lower than controls and \textbf{did not change with detoxification}. \href{https://pubmed.ncbi.nlm.nih.gov/26208874/}{PubMed} \href{https://pmc.ncbi.nlm.nih.gov/articles/PMC3261322/}{PubMed Central}

This last finding is crucial and underappreciated. Robertson et al.~(2016, \emph{Neuropsychopharmacology}) confirmed that ``there are no published reports of recovery of D2/D3 receptors with drug abstinence in human subjects.'' \href{https://pmc.ncbi.nlm.nih.gov/articles/PMC4832026/}{PubMed Central} Their small RCT (n = 19) found that 8 weeks of exercise training --- but not abstinence alone --- produced a \textbf{13.89\% increase in striatal D2/D3 availability}. \href{https://www.nature.com/articles/npp2015331}{Nature} \href{https://digitalcommons.georgefox.edu/cgi/viewcontent.cgi?article=1049&context=dmsc}{George Fox University} This remains the only published evidence that D2/D3 receptors can recover in humans after stimulant use, and it has \textbf{never been independently replicated}. The original study's limitations --- tiny sample, residential setting, no post-discharge follow-up, no control for abstinence-independent D2 changes --- make this a suggestive but unreliable finding.

For D1 receptors, which mediate a distinct set of dopamine functions, \textbf{essentially no human in vivo recovery data exists}. The glutamate and norepinephrine systems after stimulant exposure remain similarly unstudied in humans. MDMA's serotonergic effects show partial subcortical SERT recovery over months to years, but neocortical SERT shows \textbf{little evidence of recovery}, \href{https://pmc.ncbi.nlm.nih.gov/articles/PMC3521982/}{PubMed Central} and primate data reveals abnormal reinnervation patterns persisting 7 years post-exposure (Hatzidimitriou et al., 1999).

The longest follow-up PET study of methamphetamine users caps at roughly 14 months. No human longitudinal neuroimaging study has tracked stimulant recovery beyond approximately 2 years. The 4-year primate data from Woolverton et al.~(1989) --- showing significant dopamine and serotonin reductions still present in rhesus monkeys --- suggests that some deficits may never fully resolve. \href{https://pmc.ncbi.nlm.nih.gov/articles/PMC3261322/}{PubMed Central +2}

\begin{center}\rule{0.5\linewidth}{0.5pt}\end{center}

\subsection{Structural and functional brain recovery follows region-specific, heterogeneous timelines}\label{14_what_we_dont_know-structural-and-functional-brain-recovery-follows-region-specific-heterogeneous-timelines}

Structural MRI studies paint a more nuanced picture than the binary ``damaged/healed'' narrative. A large cross-sectional study (n = 99 methamphetamine users, 86 controls, \emph{BMC Psychiatry}, 2020) \href{https://bmcpsychiatry.biomedcentral.com/articles/10.1186/s12888-020-02567-3}{BioMed Central} found that hippocampal volume and insular thickness correlated positively with abstinence duration, suggesting ongoing recovery. But bilateral superior frontal gyri showed \textbf{permanently greater thickness}, attributed to gliosis \href{https://pmc.ncbi.nlm.nih.gov/articles/PMC7146984/}{PubMed Central} --- the brain's scarring response. Recovery appears region-specific: \href{https://link.springer.com/article/10.1186/s12888-020-02567-3}{Springer} cerebellar volume increases with abstinence duration; some cortical regions show rapid gray matter increases within the first month; and certain changes appear irreversible.

Longitudinal structural data confirms this heterogeneity. \href{https://www.frontiersin.org/journals/psychiatry/articles/10.3389/fpsyt.2018.00722/full}{Frontiers} Ruan et al.~tracked 30 users at 6 and 12 months of abstinence, finding that gray matter volume increased in the cerebellum while simultaneously decreasing in the cingulate gyrus \href{https://pubmed.ncbi.nlm.nih.gov/28887180/}{PubMed} --- demonstrating \textbf{dynamic, region-specific reorganization} rather than uniform healing. Kim et al.~(2006) found that long-term abstinence (≥6 months) was associated with increased gray matter density in the right middle frontal gyrus, correlated with improved Wisconsin Card Sort Test performance. \href{https://www.medrxiv.org/content/10.1101/2025.02.11.25322099v1.full}{medRxiv}

Functional recovery data is sparser. Salo et al.~(2009) found that methamphetamine users abstinent for more than 1 year showed no statistical difference from controls on the Stroop attention test \href{https://www.promises.com/addiction-blog/brain-function-can-recover-after-a-year-of-abstinence-from-methamphetamine/}{Promises Behavioral Health} \href{https://scienceblog.com/22706/attention-recovering-meth-heads-impulse-control-returns-after-a-year-gut-it-out-until-then/}{ScienceBlog} --- a finding frequently cited to support the ``12--18 month'' timeline. But this was cross-sectional, not longitudinal, and measured a single dimension of cognitive control rather than subjective experience. White matter integrity recovery, neuroinflammatory resolution, and functional connectivity normalization all remain inadequately characterized, with no true longitudinal DTI study \href{https://pubmed.ncbi.nlm.nih.gov/31176066/}{PubMed} and limited TSPO-PET data on microglial activation trajectories.

\begin{center}\rule{0.5\linewidth}{0.5pt}\end{center}

\subsection{Cognitive recovery is real but the hardest questions remain unasked}\label{14_what_we_dont_know-cognitive-recovery-is-real-but-the-hardest-questions-remain-unasked}

Meta-analytic evidence establishes that methamphetamine users show moderate cognitive deficits across multiple domains, with episodic memory most affected, followed by executive function and processing speed (Scott et al., 2007; Potvin et al., 2018). The critical question is whether these deficits reverse.

The strongest evidence comes from Iudicello et al.~(2010): 83 methamphetamine-dependent individuals assessed at baseline and 1-year follow-up, of whom 25 maintained abstinence. \href{https://pmc.ncbi.nlm.nih.gov/articles/PMC2911490/}{PubMed Central} \textbf{Abstainers showed neuropsychological performance comparable to healthy controls at 1 year}, with the most impaired at baseline showing the greatest improvement, particularly in processing speed and motor abilities. Affective distress also normalized. \href{https://pubmed.ncbi.nlm.nih.gov/20198527/}{PubMed} \href{https://pmc.ncbi.nlm.nih.gov/articles/PMC2911490/}{PubMed Central} Basterfield et al.'s 2019 meta-analysis of 1,008 abstinent methamphetamine users found only small-to-moderate persistent effects, with \textbf{three domains --- fluid reasoning, short-term working memory, and reaction speed --- showing no group differences} from controls at all. They cautioned that ``strong statements regarding impaired cognitive functioning in abstinent methamphetamine users are premature.'' \href{https://unitedstatespsychiatricassociation.com/general/a-meta-analysis-of-the-relationship-between-abstinence-and-neuropsychological-functioning-in-methamphetamine-use-disorder/}{Unitedstatespsychiatricassociation}

Processing speed appears to be the most recovery-responsive domain, with substantial improvement within 6--12 months. Executive function shows meaningful recovery by 6--12 months. Memory remains the most persistently impaired domain, with inconsistent evidence for full normalization. One troubling finding from Huckans et al.~(2021) showed that individuals in remission from methamphetamine performed worse than controls on memory indices \href{https://pubmed.ncbi.nlm.nih.gov/34612792/}{PubMed} --- potentially reflecting either persistent damage or pre-existing vulnerability that only manifests clearly in the absence of stimulant-enhanced performance.

Yet here the measurement problem becomes acute. Cognitive tests measure whether someone can sort cards or suppress a prepotent response in a laboratory. They do not measure whether someone can enjoy a sunset, feel motivated to pursue goals, or experience genuine connection with another person. \textbf{The ecological validity of cognitive improvement --- whether normalized lab performance translates to real-world functional capacity --- is almost never assessed.}

\begin{center}\rule{0.5\linewidth}{0.5pt}\end{center}

\subsection{Anhedonia is the most important outcome that nobody measures}\label{14_what_we_dont_know-anhedonia-is-the-most-important-outcome-that-nobody-measures}

Wang et al.'s 2004 finding of persistent nucleus accumbens hypometabolism is the closest the field has come to a neurobiological explanation for protracted anhedonia. \href{https://ohsu.elsevierpure.com/en/publications/partial-recovery-of-brain-metabolism-in-methamphetamine-abusers-a-2/}{Elsevierpure} Clinical descriptions universally acknowledge it as perhaps the most debilitating and persistent symptom of methamphetamine recovery. Treatment center literature commonly cites timelines of ``6--12 months, up to 2 years'' for hedonic recovery. SAMHSA clinical guidance notes that ``it may take 6--12 months for normal pleasure response to return.'' \href{https://healingpinesrecovery.com/stimulant-withdrawal/}{Healingpinesrecovery}

\textbf{None of these claims are supported by systematic empirical measurement.} No published longitudinal study has tracked anhedonia severity over time in methamphetamine recovery using validated anhedonia scales such as the SHAPS, TEPS, or DARS. No validated instrument exists specifically for measuring the temporal course of post-stimulant anhedonia resolution. The general anhedonia scales were designed for depression, not for the potentially distinct neurobiological substrate of post-stimulant reward system dysfunction.

The same void applies to amotivation. Clinical reports of ``lack of spontaneity'' --- a flatness of drive and initiative that family members describe in their recovering loved ones --- appear nowhere in the quantitative literature. Energy, motivation, sleep architecture recovery, personality changes noted by family members, and the capacity for emotional depth are all clinically salient and completely unmeasured.

A 2025 study in \emph{Harm Reduction Journal} finally asked 100 methamphetamine users what recovery outcomes matter to them. \href{https://pmc.ncbi.nlm.nih.gov/articles/PMC11740499/}{PubMed Central} \href{https://harmreductionjournal.biomedcentral.com/articles/10.1186/s12954-025-01155-6}{Harm Reduction Journal} They prioritized: preventing legal trouble (92\%), employment stability (82\%), \textbf{improved quality of life (80\%)}, housing stability (78\%), improved coping (78\%), improved relationships (75\%). \href{https://pubmed.ncbi.nlm.nih.gov/39825407/}{PubMed} \href{https://pmc.ncbi.nlm.nih.gov/articles/PMC11740499/}{PubMed Central} These are rarely primary endpoints in clinical trials, which overwhelmingly measure binary abstinence status and days of use.

The only longitudinal quality-of-life study specifically in methamphetamine users --- Gonzales et al.~(2009), using the SF-8 in 723 participants --- found that mental health-related quality of life improved as a function of treatment completion, while physical health showed ``fairly static trajectories.'' No study has measured quality of life in former methamphetamine users at 2, 5, or 10 years post-cessation. The temporal gap table is stark:

{\def\LTcaptype{none} % do not increment counter
\begin{longtable}[]{@{}
  >{\raggedright\arraybackslash}p{(\linewidth - 6\tabcolsep) * \real{0.1594}}
  >{\raggedright\arraybackslash}p{(\linewidth - 6\tabcolsep) * \real{0.2174}}
  >{\raggedright\arraybackslash}p{(\linewidth - 6\tabcolsep) * \real{0.3043}}
  >{\raggedright\arraybackslash}p{(\linewidth - 6\tabcolsep) * \real{0.3188}}@{}}
\toprule\noalign{}
\begin{minipage}[b]{\linewidth}\raggedright
Timescale
\end{minipage} & \begin{minipage}[b]{\linewidth}\raggedright
Anhedonia data
\end{minipage} & \begin{minipage}[b]{\linewidth}\raggedright
Quality of life data
\end{minipage} & \begin{minipage}[b]{\linewidth}\raggedright
Phenomenological data
\end{minipage} \\
\midrule\noalign{}
\endhead
\bottomrule\noalign{}
\endlastfoot
0--6 months & Described clinically, never measured & None & None \\
6--12 months & Inferred from Wang 2004 PET only & One study (Gonzales 2009) & None \\
1--2 years & None & None & None \\
2--5 years & \textbf{None} & \textbf{None} & \textbf{None} \\
5--10 years & \textbf{None} & \textbf{None} & \textbf{None} \\
10+ years & \textbf{None} & \textbf{None} & \textbf{None} \\
\end{longtable}
}

\begin{center}\rule{0.5\linewidth}{0.5pt}\end{center}

\subsection{Clinical outcomes reveal severe mortality risk and a treatment vacuum}\label{14_what_we_dont_know-clinical-outcomes-reveal-severe-mortality-risk-and-a-treatment-vacuum}

The Taiwanese National Health Insurance Research Database has produced the most rigorous mortality data. Lee et al.~(2021) followed \textbf{23,248 individuals} with methamphetamine use disorder from 2001--2005 and found an all-cause standardized mortality ratio (SMR) of \textbf{5.4} --- meaning methamphetamine users died at 5.4 times the rate of the general population. Unnatural cause SMR was 14.8. \href{https://onlinelibrary.wiley.com/doi/abs/10.1111/add.15501}{Wiley Online Library} \textbf{Suicide had the highest SMR at 16.3}, with drug overdose suicide at 24.9. \href{https://www.researchgate.net/publication/41423198_Three-year_mortality_and_predictors_after_release_A_longitudinal_study_of_the_first-time_drug_offenders_in_Taiwan}{ResearchGate} Mean age of death was 38.4 years --- a loss of more than 30 years of potential life. \href{https://pubmed.ncbi.nlm.nih.gov/38462530/}{PubMed} \href{https://www.jstage.jst.go.jp/article/jea/34/10/34_JE20230263/_article}{JST} Kuo et al.~(2011) confirmed these findings in a separate cohort of 1,254 patients \href{https://pmc.ncbi.nlm.nih.gov/articles/PMC3117111/}{PubMed Central} (SMR = 6.02), with women showing dramatically higher unnatural death ratios \href{https://pmc.ncbi.nlm.nih.gov/articles/PMC3117111/}{PubMed Central} (SMR 26.19 vs.~9.82 for men). \href{https://pmc.ncbi.nlm.nih.gov/articles/PMC3117111/}{PubMed Central}

Treatment outcomes are discouraging. No FDA-approved pharmacotherapy exists. \href{https://nida.nih.gov/publications/drugs-brains-behavior-science-addiction/treatment-recovery}{NIDA} A systematic review of 43 RCTs testing 23 pharmacotherapies (Siefried et al., 2020) found \textbf{no convincing results}, with the most promising signals from stimulant agonists (dexamphetamine, methylphenidate), naltrexone, and topiramate --- all rated as weak or inconsistent evidence. Behavioral treatments fare somewhat better: contingency management shows the strongest evidence, \href{https://www.sciencedirect.com/science/article/abs/pii/S009174352300230X}{ScienceDirect +2} and the Matrix Model produces measurable improvements, but long-term follow-up is limited. The Calabria et al.~(2010) systematic review found only \textbf{one} prospective study with ≥3-year follow-up that met inclusion criteria for amphetamine dependence remission, \href{https://www.sciencedirect.com/science/article/abs/pii/S0306460310000985}{ScienceDirect} estimating a conservative rate of 16\% per year. \href{https://findings.org.uk/PHP/dl.php?f=Calabria_B_1.cab&s=dd&sf=mx}{Findings}

The longest published follow-up of methamphetamine users in the English-language literature is the Hser combined cohort analysis covering approximately 474 methamphetamine users with a mean follow-up of \textbf{20.9 years}. \href{https://pmc.ncbi.nlm.nih.gov/articles/PMC4072125/}{PubMed Central} But this tracked drug use patterns --- daily use, non-daily use, abstinence, incarceration --- not neurobiological recovery, quality of life, or subjective well-being. \href{https://pmc.ncbi.nlm.nih.gov/articles/PMC4072125/}{PubMed Central} We know something about whether people continue using over two decades. We know almost nothing about what their lives are like.

\begin{center}\rule{0.5\linewidth}{0.5pt}\end{center}

\subsection{Animal models promise translation they cannot deliver}\label{14_what_we_dont_know-animal-models-promise-translation-they-cannot-deliver}

The translational architecture of amphetamine neurotoxicity research is built on a foundation of rodent data that demonstrably fails to predict primate --- let alone human --- recovery. Recovery speed is \textbf{inversely related to phylogenetic proximity to humans}: mice recover dopaminergic function in days to weeks, rats in months, \href{https://pmc.ncbi.nlm.nih.gov/articles/PMC3261322/}{PubMed Central} vervet monkeys over months to years, and rhesus monkeys show deficits persisting 4 years. \href{https://www.jneurosci.org/content/21/23/9414}{Journal of Neuroscience}

The Melega lab's work is pivotal. Melega et al.~(1997) demonstrated that vervet monkeys given acute neurotoxic methamphetamine doses showed FDOPA Ki decreases of 60--70\% at 1--2 weeks, recovering to only a 16\% deficit at 32 weeks \href{https://www.sciencedirect.com/science/article/abs/pii/S0006899397005489}{ScienceDirect} --- a timeline orders of magnitude slower than comparable mouse studies. Melega et al.~(2008) created the most human-relevant primate model: escalating doses over 33 weeks in socially housed vervets, targeting plasma methamphetamine concentrations of 1--3 μM (matching human abuser levels). \href{https://www.nature.com/articles/1301502}{Nature} Even this pharmacologically conservative paradigm produced \textbf{20\% lower striatal dopamine and 35\% lower DAT binding} at 3 weeks post-cessation. \href{https://www.researchgate.net/publication/6211685_Long-Term_Methamphetamine_Administration_in_the_Vervet_Monkey_Models_Aspects_of_a_Human_Exposure_Brain_Neurotoxicity_and_Behavioral_Profiles}{ResearchGate} Behavioral differences normalized within 3 weeks, but neurochemical damage persisted --- precisely mirroring the human dissociation between superficial functional recovery and underlying neurochemical deficit.

The Woolverton et al.~(1989) rhesus monkey study remains the longest primate follow-up: \textbf{significant dopamine and serotonin reductions were still evident 4 years after methamphetamine treatment}. \href{https://pmc.ncbi.nlm.nih.gov/articles/PMC8338805/}{PubMed Central} No cell body loss was observed, distinguishing the primate pathology from mouse models where substantia nigra cell death occurs --- a species-specific vulnerability that does not appear to generalize. \href{https://www.jneurosci.org/content/18/20/8417}{Journal of Neuroscience}

Groman et al.~(2012) demonstrated that chronic escalating methamphetamine in vervets produced D2 receptor decreases and reversal learning deficits \href{https://www.jneurosci.org/content/32/17/5843}{Journal of Neuroscience} \href{https://www.ncbi.nlm.nih.gov/pmc/articles/PMC3353813/}{NCBI} persisting beyond 7 weeks, with marked individual heterogeneity. Villemagne et al.~(1998) showed DAT neurotoxicity in baboons at all doses tested, including 0.5 mg/kg × 4 --- which scales to approximately \textbf{26 mg for a 70 kg human}, well within a single recreational dose. \href{https://pmc.ncbi.nlm.nih.gov/articles/PMC6793413/}{PubMed Central}

The C57BL/6 mouse --- the workhorse of the field --- is a particularly poor behavioral model. Grace et al.~(2010) demonstrated that mice given neurotoxic methamphetamine doses producing dopamine depletion comparable to rats showed \textbf{no learning or memory deficits} on standard tests, concluding that this ``may limit the utility of C57BL/6 mice for modeling the cognitive and behavioral effects described in human MA users.'' \href{https://pmc.ncbi.nlm.nih.gov/articles/PMC2854319/}{PubMed Central} The field's most commonly used model organism cannot reproduce the cognitive effects that are among the most clinically significant consequences of the drug.

Dose equivalence problems compound the translation failure. The standard mouse binge paradigm (4 × 10 mg/kg at 2-hour intervals = 40 mg/kg/day) converts to a human equivalent dose of approximately 227 mg/day using FDA-recommended body surface area scaling \href{https://pmc.ncbi.nlm.nih.gov/articles/PMC4804402/}{PubMed Central} --- the high end of heavy use. But pharmacokinetic differences are enormous: methamphetamine half-life is approximately 1 hour in mice versus 10--12 hours in humans, meaning the 2-hour dosing interval creates an entirely different exposure profile. Melega's approach of matching plasma concentrations rather than mg/kg doses represents a far superior methodology, but it requires primate subjects at 100--500× the cost.

Species that might bridge the rodent--primate gap remain almost entirely unstudied. \textbf{Pigs, whose dopamine systems and brain architecture closely resemble those of humans, have virtually zero publications on methamphetamine recovery.} No dog, ferret, or New World monkey studies exist. Zebrafish offer a unique advantage --- as poikilotherms, they eliminate the hyperthermia confound that complicates all mammalian neurotoxicity data \href{https://www.frontiersin.org/journals/pharmacology/articles/10.3389/fphar.2022.976970/full}{Frontiers} --- but the first zebrafish methamphetamine neurotoxicity model was published only in 2021 (Bedrossiantz et al.) \href{https://www.frontiersin.org/journals/pharmacology/articles/10.3389/fphar.2021.770319/full}{Frontiers} and no recovery data exists. C. elegans \href{https://pubmed.ncbi.nlm.nih.gov/26810004/}{PubMed} and Drosophila are used for genetic mechanism discovery \href{https://www.mdpi.com/2227-9059/10/1/119}{MDPI} but cannot model recovery trajectories. The phylogenetic sampling is so sparse that \textbf{we cannot determine whether the mouse-to-primate recovery difference represents a gradual phylogenetic gradient, a primate-specific trait, or a body-size scaling phenomenon}.

The publication landscape reflects the structural problem. For every primate methamphetamine paper, approximately \textbf{20--40 rodent papers} exist. The cost differential is the primary driver: \href{https://www.ncbi.nlm.nih.gov/books/NBK100126/}{NCBI} a basic mouse R01 fits within NIH's \$250,000/year modular budget; \href{https://www.ninds.nih.gov/funding/preparing-your-application/preparing-budget/creating-budget}{National Institute of Neurological Disorders and Stroke} \href{https://biosciencewriters.com/Should-I-Apply-for-an-R01-R03-or-R21-NIH-Grant-Award.aspx}{Biosciencewriters} a longitudinal primate PET study easily requires \$500,000--\$2,000,000 in direct costs. A mouse experiment produces publishable data within weeks; a primate study requires months to years. Junior faculty facing tenure review at year 6 cannot afford to start a decade-long primate study that may produce only 1--2 papers.

\begin{center}\rule{0.5\linewidth}{0.5pt}\end{center}

\subsection{History's largest natural experiments remain inaccessible}\label{14_what_we_dont_know-historys-largest-natural-experiments-remain-inaccessible}

Three historical episodes generated methamphetamine-exposed populations on a scale that dwarfs any clinical study, yet the follow-up data is either lost, language-locked, or was never collected.

Japan's postwar Hiropon epidemic \href{https://www.cambridge.org/core/journals/journal-of-asian-studies/article/abs/methamphetamine-solution-drugs-and-the-reconstruction-of-nation-in-postwar-japan/0A0FCE7E74AE7D1ECE9A134114773CE6}{Cambridge Core} exposed an estimated \textbf{1.5--2 million people} to methamphetamine between 1945 and the mid-1950s. \href{https://www.unodc.org/unodc/en/data-and-analysis/bulletin/bulletin_1989-01-01_1_page007.html}{UNODC} \href{https://www.zocalopublicsquare.org/the-world-war-ii-wonder-drug-that-never-left-japan/}{Zócalo Public Square} Tatetsu, Goto, and Fujiwara published the foundational monograph \emph{Kakuseizai Chūdoku} (The Methamphetamine Psychosis) in 1956 --- the first systematic clinical description of methamphetamine psychosis, from the pre-neuroleptic era. \textbf{No English translation exists.} It is known almost exclusively through secondary citations in Sato and Ujike's later English-language papers. Tatetsu reportedly conducted a 30-year follow-up of 9 patients circa 1997; this publication could not be located in any indexed database and appears to exist, if at all, only in Japanese psychiatric archives. Teraoka's 1967 8--12 year follow-up is similarly lost.

What happened to 2 million users is the largest unanswered question in the field's history. No accessible epidemiological follow-up exists. The epidemic was suppressed through legislation and policing, \href{https://journals.sagepub.com/doi/10.1177/009145090803500410}{Sage Journals} \href{https://www.isdp.eu/publication/forgotten-success-story-japan-methamphetamine-problem/}{Institute for Security and Development Policy} but the clinical fate of this population --- whether most recovered, how many developed chronic psychosis, what their long-term functional outcomes were --- remains unknown. Sato et al.~(1983) established that psychotic relapse can occur after years of complete abstinence, triggered by stress \href{https://nyaspubs.onlinelibrary.wiley.com/doi/10.1111/j.1749-6632.2000.tb05178.x}{Annals of the New York Academy of Sciences} --- a finding of enormous clinical significance that emerged from decades of Japanese case observation but whose full empirical basis remains locked in Japanese-language publications.

The German Pervitin literature is similarly inaccessible. Over 100 clinical papers on Pervitin (methamphetamine) were published by late 1939 in German medical journals. \href{https://scholarworks.uttyler.edu/cgi/viewcontent.cgi?article=1001&context=senior_projects}{Uttyler} Staehelin's 1941 ``Pervitin-Psychose'' paper from Basel, Greving's 1941 report on ``psychopathological and physical processes in years-long Pervitin abuse,'' \href{https://www.researchgate.net/publication/275135440_Speed_in_the_Third_Reich_Metamphetamine_Pervitin_Use_and_a_Drug_History_From_Below}{ResearchGate} and numerous other clinical descriptions exist behind Springer paywalls and in German-language physical archives. Millions of Wehrmacht soldiers were exposed to Pervitin; \href{https://en.wikipedia.org/wiki/Drug_policy_of_Nazi_Germany}{Wikipedia} their postwar psychiatric outcomes would be extraordinarily valuable but appear never to have been systematically studied, or such studies exist only in inaccessible German archives.

The Czech Republic's 50-year methamphetamine (pervitin) tradition \href{https://journals.sagepub.com/doi/10.1177/002204260703700108}{Sage Journals} \href{https://www.researchgate.net/publication/354534700_Methamphetamine_use_and_consequences_in_context_of_drug_situation_in_the_Czech_Republic}{ResearchGate} has produced important but largely Czech-language research. Nepustil's qualitative study of 19 people who quit methamphetamine for more than 5 years without professional treatment \href{https://www.academia.edu/35726007/Recovered_without_Treatment_The_Process_of_Abandoning_Crystal_Meth_Use_without_Professional_Help_Pavel_Nepustil_A_Taos_Institute_Publication}{Academia.edu} (published in English in 2013) represents some of the only data on long-term natural recovery --- finding that the process centered on disruption of relational patterns, co-creation of transitional trajectories, and development of new belonging rather than on neurobiological events. \href{https://www.taosinstitute.net/product/recovered-without-treatment-the-process-of-abandoning-crystal-meth-use-without-professional-help-by-pavel-nepustil}{The Taos Institute} Šefránek et al.~(2017) published therapeutic community outcomes in English. \href{https://www.tandfonline.com/doi/full/10.1080/1556035X.2017.1280718}{Taylor \& Francis Online} But the communist-era research (including Večerková 1986) appears irretrievably lost.

China's current methamphetamine epidemic --- with registered users exceeding 1 million and a National Dynamic Management Database containing data on nearly 2 million cases \href{https://pmc.ncbi.nlm.nih.gov/articles/PMC9096246/}{PubMed Central} --- represents perhaps the world's largest potential cohort for recovery research. But this data remains \textbf{language-locked behind CNKI (China National Knowledge Infrastructure) paywalls and Chinese-language barriers}. No accessible English-language publication of an 8-year or longer follow-up Chinese methamphetamine cohort was identified.

\begin{center}\rule{0.5\linewidth}{0.5pt}\end{center}

\subsection{Why the gaps exist: a structural analysis of research neglect}\label{14_what_we_dont_know-why-the-gaps-exist-a-structural-analysis-of-research-neglect}

The absence of long-term phenomenological recovery data is not a random oversight. It emerges from the convergence of at least five structural forces.

\textbf{Funding architecture favors fast, controllable science.} NIDA's annual budget of approximately \textbf{\$1.6 billion} funds over 85\% of global addiction research. \href{https://en.wikipedia.org/wiki/National_Institute_on_Drug_Abuse}{Wikipedia} A modular R01 grant (\textasciitilde\$250,000/year for 5 years) \href{https://www.ninds.nih.gov/funding/preparing-your-application/preparing-budget/creating-budget}{National Institute of Neurological Disorders and Stroke} neatly supports a mouse study with one postdoc, one graduate student, and a colony of 150 cages (\textasciitilde\$42,000/year). \href{https://drugmonkey.scientopia.org/2013/03/11/a-full-modular-250k-per-year-nih-grant-doesnt-actually-pay-for-itself/}{Drugmonkey} A longitudinal human PET study with 50 participants scanned at 4 timepoints over 5 years requires \$500,000--\$1,000,000+ in direct costs --- exceeding modular limits, requiring detailed budget justification, and triggering greater reviewer skepticism. Congress earmarked only \textbf{\$10 million} specifically for stimulant use disorder treatment research in FY2021, \href{https://nasadad.org/wp-content/uploads/2020/11/FY-2021-Senate-Recommendations.pdf}{Nasadad} against a backdrop of hundreds of millions for opioid research.

\textbf{Career incentives systematically favor mouse studies.} Mouse experiments generate publishable data in weeks; each variation (dose, timepoint, brain region, strain) becomes a separate paper. A productive mouse lab on a single R01 can produce 3+ papers per grant cycle --- sufficient for tenure. A 10-year human longitudinal study started as an assistant professor would yield no data by tenure review. NIH study sections reviewing addiction grants are populated by basic scientists trained in animal models, creating self-reinforcing funding preferences. The tools of mouse neuroscience --- optogenetics, CRISPR, fiber photometry --- are considered cutting-edge and publishable in high-impact journals; patient-reported outcome measurement is not.

\textbf{The measurement paradigm equates abstinence with recovery.} The field's dominant outcome measure is binary use status, supplemented by urine drug screens and days of use. The FDA required complete abstinence as the endpoint for SUD medication approval until recent shifts. \href{https://harmreductionjournal.biomedcentral.com/articles/10.1186/s12954-025-01155-6}{Harm Reduction Journal} By defining recovery as abstinence, the field excludes examination of the quality of abstinent life --- whether it is characterized by genuine well-being or by endured misery. Paquette et al.~(2022) documented that nearly 100,000 U.S. treatment episodes annually end in administrative discharge for any substance use during treatment, embedding the abstinence framework in institutional practice. \href{https://pmc.ncbi.nlm.nih.gov/articles/PMC8815796/}{PubMed Central}

\textbf{Stigma shapes which questions are considered worth asking.} Methamphetamine users face greater stigmatizing attitudes than those with cocaine, opioid, or alcohol use disorders (ETSU dissertation study). Greater perceived moral blameworthiness correlates with reduced investment in studying recovery --- if the person caused their own condition, the implicit logic goes, their subjective experience of recovery is less worthy of scientific attention. Healthcare providers report treating methamphetamine users as making ``poor choices,'' \href{https://pmc.ncbi.nlm.nih.gov/articles/PMC10896942/}{PubMed Central} and this orientation extends to research design: the field asks ``did they stop using?'' rather than ``are they able to live a full life?''

\textbf{The translational pipeline produces knowledge without utility.} Venniro et al.~(2020, \emph{Nature Reviews Neuroscience}) stated plainly that animal model ``mechanistic gains have not led to improvements in addiction treatment.'' \href{https://www.nature.com/articles/s41583-020-0378-z}{Nature} Despite thousands of rodent studies identifying molecular targets --- VMAT2, TAAR1, various kinases, inflammatory pathways --- no FDA-approved medication for methamphetamine use disorder has emerged. \href{https://nida.nih.gov/about-nida/legislative-activities/budget-information/fiscal-year-2022-budget-information-congressional-justification-national-institute-drug-abuse}{NIDA} The most promising clinical result, bupropion combined with naltrexone from the ADAPT-2 trial, was found through clinical observation rather than rodent-to-human translation. \href{https://nida.nih.gov/about-nida/legislative-activities/budget-information/fiscal-year-2022-budget-information-congressional-justification-national-institute-drug-abuse}{NIDA} The field has created what one might call a knowledge production system optimized for publications rather than clinical impact.

\begin{center}\rule{0.5\linewidth}{0.5pt}\end{center}

\subsection{The questions that have never been asked}\label{14_what_we_dont_know-the-questions-that-have-never-been-asked}

The most damning finding of this review is not what the literature shows but what it has never attempted to measure. The following questions --- all of fundamental clinical importance --- have \textbf{no systematic empirical data} supporting any answer:

\begin{itemize}
\tightlist
\item
  When does hedonic capacity fully return after methamphetamine cessation? (No data beyond \textasciitilde14 months)
\item
  Is there a threshold of use duration or intensity beyond which full hedonic recovery is impossible?
\item
  What is the natural trajectory of anhedonia resolution from month 6 to year 10?
\item
  Do former methamphetamine users ever return to pre-use baseline on measures of pleasure, motivation, and meaning?
\item
  What predicts who recovers hedonic function and who does not --- genetics, duration of use, age of onset, social support?
\item
  How does protracted post-stimulant anhedonia differ from clinical depression, and does it require different treatment?
\item
  What do people who have maintained 5+ years of methamphetamine abstinence report about their subjective experience?
\end{itemize}

The longest published human neuroimaging follow-up is approximately \textbf{14 months}. The longest behavioral follow-up (Hser combined cohorts, \textasciitilde20.9 years) measures use patterns, not well-being. \href{https://pmc.ncbi.nlm.nih.gov/articles/PMC4072125/}{PubMed Central} \href{https://pubmed.ncbi.nlm.nih.gov/24837584/}{PubMed} There are zero validated instruments designed for tracking post-stimulant anhedonia recovery. There are zero systematic studies of the phenomenology of methamphetamine recovery beyond 5 years. The period between years 2 and 10 --- precisely when sustained recovery is consolidated or fails --- is a \textbf{complete evidential void}.

\begin{center}\rule{0.5\linewidth}{0.5pt}\end{center}

\subsection{Conclusion: the scope of what we do not know is itself a finding}\label{14_what_we_dont_know-conclusion-the-scope-of-what-we-do-not-know-is-itself-a-finding}

This review documents a research landscape defined by a paradox: extraordinary molecular detail about how methamphetamine damages dopamine terminals in C57BL/6 mice coexists with near-total ignorance about whether human beings who stop using the drug can eventually enjoy their lives again. The foundational ``brain healing'' narrative rests on 5 longitudinal PET subjects. D2/D3 receptors --- arguably the most functionally relevant dopamine biomarker for subjective experience --- show no spontaneous recovery in humans. \href{https://pmc.ncbi.nlm.nih.gov/articles/PMC4832026/}{PubMed Central} The most commonly used animal model cannot reproduce the drug's cognitive effects in humans, \href{https://pmc.ncbi.nlm.nih.gov/articles/PMC2854319/}{PubMed Central} recovers on timescales that dramatically underestimate the human experience, and has failed to produce a single approved medication. History's largest natural experiments --- 2 million Japanese Hiropon users, \href{https://www.zocalopublicsquare.org/the-world-war-ii-wonder-drug-that-never-left-japan/}{Zócalo Public Square} millions of German soldiers exposed to Pervitin \href{https://en.wikipedia.org/wiki/Drug_policy_of_Nazi_Germany}{Wikipedia} --- generated clinical observations that are language-locked, paywalled, or lost.

What would need to change is not mysterious. A dedicated NIDA Recovery Science Initiative with ring-fenced funding for prospective studies exceeding 5 years. \href{https://www.science.org/doi/10.1126/sciadv.aaz6596}{Science} A validated instrument for post-stimulant anhedonia and hedonic recovery. Protected career pathways for researchers who choose slow, difficult human work over fast, publishable mouse experiments. Mandatory patient-reported outcomes as co-primary endpoints in every treatment trial. Expansion of species diversity in preclinical research --- particularly pigs, whose dopamine systems most closely resemble those of humans among non-primates. And a fundamental reorientation of the field's central question from ``did they stop using?'' to ``can they live a life worth living?''

The gap between what we know about methamphetamine damage and what we know about methamphetamine recovery is not merely a gap. It is a statement about whose suffering the scientific enterprise considers worth understanding.
